\section{Computational cognitive models: formal boundaries}
\label{sec:p19-ccm-annex}

The tensor results make precise the equations at printed 60--61;
the network result qualifies the overbroad capacity statement at 57
\parencite{DoumasHummel2012HigherCognition}. These proofs do not validate
a learning mechanism, psychological interpretation or neural implementation.

\begin{annexitem}{annex:p19.ccm.continuity}{thm:p19-ccm-continuity}{Finite continuous layers and exact rational-input steps}
\begin{proof}
Each coordinate of an affine layer is a constant plus a finite sum
of scalar multiples of input coordinates, hence is continuous.
Coordinatewise application of continuous scalar activations is
continuous. Composing finitely many such layers therefore gives a
continuous map; this also covers constant layers and zero weights.

Suppose a scalar output realizes the asserted rational step.
At input zero its value is zero. For every positive integer \(n\),
input \(1/n\) is rational and positive, so its value is one.
But \(1/n\) converges to zero, and continuity requires those
output values to converge to the value at zero, namely zero.
The constant sequence of ones does not converge to zero, a
contradiction. The step's rational-input computational rule is
decidable: compare the numerator of a finite rational encoding with
zero, after using a positive denominator. This does not supply a
continuous extension of that step-valued rule.
\end{proof}
\end{annexitem}

\begin{annexitem}{annex:p19.ccm.tensor-similarity}{thm:p19-ccm-tensor-similarity}{All tensor identities and their boundary cases}
\begin{proof}
Flattening retains one coordinate for every pair \((i,j)\).
Finite distributivity gives
\[
 \sum_{i,j}(r_i^{(1)}f_j^{(1)})(r_i^{(2)}f_j^{(2)})
 =\left(\sum_i r_i^{(1)}r_i^{(2)}\right)
  \left(\sum_j f_j^{(1)}f_j^{(2)}\right).
\]
This is the dot-product identity, including zero factors.
Applying it to a tensor paired with itself gives
\(\|r\otimes f\|_2^2=\|r\|_2^2\|f\|_2^2\).
Taking nonnegative square roots gives
\(\|r\otimes f\|_2=\|r\|_2\|f\|_2\).
For four nonzero factors, divide the dot identity by the two
nonzero tensor norms. With cosine defined as dot product divided
by the product of Euclidean norms, this gives the cosine identity.
If a factor is zero, its tensor is zero, and that cosine division
is undefined; the unnormalized dot identity still holds.

The three-factor case uses exactly the same finite expansion:
the sum over \((i,j,k)\) of a product separates into the three
coordinate sums. Pairing with itself then gives the product of
three squared norms, and positive norms permit the corresponding
cosine division. No infinite-sum or limiting interchange is used.

For an identical nonzero role, its self-cosine is one. Orthogonal
nonzero fillers have dot product and cosine zero, so their bound
cosine is zero. Finally every coordinate of
\((-r)\otimes(-f)\) equals \((-r_i)(-f_j)=r_i f_j\), so
this tensor is identical to \(r\otimes f\). Thus whole-bound
identity alone does not uniquely identify unrestricted factors.
\end{proof}
\end{annexitem}

\begin{annexitem}{annex:p19.ccm.unbinding}{thm:p19-ccm-unbinding}{The decoder and arbitrary novel fillers}
\begin{proof}
The outer product \(r^{(i)}\otimes f^{(i)}\) is the matrix
\(r^{(i)}(f^{(i)})^{\mathsf T}\). Transpose the finite sum and
multiply by \(r^{(j)}\):
\[
 T^{\mathsf T}r^{(j)}
 =\sum_{i=1}^k f^{(i)}(r^{(i)})^{\mathsf T}r^{(j)}
 =\sum_{i=1}^k f^{(i)}\bigl(r^{(i)}\cdot r^{(j)}\bigr)
 =f^{(j)}.
\]
Orthonormality makes the scalar one for \(i=j\) and zero
otherwise. The identity holds for arbitrary real fillers,
including zero fillers and ones absent from any previous example.
It proves the stipulated decoder's recovery, not a learning result.
For the excluded nonorthogonal case, the same first two equalities
retain cross terms, which need not vanish; hence the stated
orthogonality premise cannot be silently dropped.
\end{proof}
\end{annexitem}
